% Relay payload classes: bent-pipe transponder, regenerative payload with on-board bundle node, navigation payload.
\begin{tikzpicture}[x=1cm,y=1cm,font=\scriptsize,
  blk/.style={draw=black!70, fill=white, minimum height=0.5cm, minimum width=1.5cm, align=center, inner sep=2pt},
  reg/.style={blk, fill=blue!10}, nav/.style={blk, fill=green!12},
  ar/.style={-{Stealth[length=4pt]}, thick, black!70},
  note/.style={rounded corners=2pt, inner xsep=4pt, inner ysep=3pt, align=left, font=\scriptsize, text width=3.7cm, line width=0.5pt}]
% (a) bent pipe
\node[font=\footnotesize\bfseries,anchor=west] at (0,2.4) {(a) Bent-pipe (transparent) relay};
\node[blk] (a1) at (0.9,1.6) {Rx antenna\\S or Ka};
\node[blk] (a2) at (2.8,1.6) {LNA};
\node[blk] (a3) at (4.7,1.6) {Frequency\\converter};
\node[blk] (a4) at (6.6,1.6) {HPA\\(TWTA or SSPA)};
\node[blk] (a5) at (8.5,1.6) {Tx antenna\\Ka or X};
\draw[ar] (a1)--(a2); \draw[ar] (a2)--(a3); \draw[ar] (a3)--(a4); \draw[ar] (a4)--(a5);
\node[note, draw=red!55, fill=red!5, text=red!60!black, anchor=west] (na) at (11.6,1.6) {\textbf{No storage:} relays only when user and Earth station are in view};
% (b) regenerative
\node[font=\footnotesize\bfseries,anchor=west] at (0,0.3) {(b) Regenerative relay with on-board networking (LunaNet store-and-forward service)};
\node[blk] (b1) at (0.9,-0.6) {Rx antenna};
\node[blk] (b2) at (2.8,-0.6) {Demodulate,\\decode};
\node[reg] (b3) at (4.9,-0.6) {Bundle node (BPv7)\\LTP, CGR, BPSec};
\node[reg] (b4) at (4.9,-1.7) {Storage\\(far-side S\&F)};
\node[blk] (b5) at (7.2,-0.6) {Encode,\\modulate};
\node[blk] (b6) at (8.9,-0.6) {HPA};
\node[blk] (b7) at (10.5,-0.6) {Tx antenna};
\draw[ar] (b1)--(b2); \draw[ar] (b2)--(b3); \draw[ar] (b3)--(b5); \draw[ar] (b5)--(b6); \draw[ar] (b6)--(b7); \draw[ar,<->] (b3)--(b4);
\node[note, draw=blue!45, fill=blue!5, text=black!75, anchor=west] (nb) at (11.6,-0.6) {\textbf{Gain:} per-hop coding gain and routing; \textbf{Cost:} router mass, power and radiation-tolerant compute};
% (c) navigation payload sharing the frequency reference
\node[font=\footnotesize\bfseries,anchor=west] at (0,-2.7) {(c) Navigation payload};
\node[nav] (c1) at (1.3,-3.6) {GNSS receiver\\(Earth GNSS, weak signal)};
\node[nav] (c2) at (4.0,-3.6) {Ultra-stable\\oscillator};
\node[nav] (c3) at (6.6,-3.6) {Navigation signal\\generator (AFS, S-band)};
\node[nav] (c4) at (9.3,-3.6) {Wide-beam\\S-band antenna};
\draw[ar] (c1)--(c2); \draw[ar] (c2)--(c3); \draw[ar] (c3)--(c4);
\draw[ar,dashed] ($(c2.north)+(-0.4,0)$) |- ($(b2.south)+(0,-0.45)$) -- (b2.south);
\node[anchor=west,text=black!60,font=\scriptsize\itshape] at ($(c2.north)+(-0.3,0.95)$) {shared frequency reference};
\node[note, draw=green!45!black, fill=green!5, text=black!75, anchor=west] (nc) at (11.6,-3.6) {Time steered to GPS time between fixes; ranging and data share \textbf{one power budget}};
\end{tikzpicture}
